前三部分讨论了构建先进大语言模型智能体所需的人工智能技术，使这些智能体能够与现实世界交互并采取行动。凭借这些能力，智能体可以为用户完成有用的工作。本部分讨论大语言模型如何使用工具执行专门任务，例如查找信息或控制机器人。智能体通过学会使用应用程序编程接口（Application Programming Interface，API）来调用工具。

工具能够帮助我们解决预训练大语言模型的一项主要局限：模型内部信息的新近程度受其训练数据集所限。检索增强技术使大语言模型能够提供最新信息。

人们已经构建了许多此类智能体。本部分将回顾医学、金融和科学等领域的重要实例，最后通过一个实践案例把这些技术结合起来，实现一个用于医学诊断的大语言模型智能体。
